\documentclass[11pt,a4paper]{article}
% =====================================================================
%  Shared preamble for the Part V lecture notes (lectures 19-22).
%
%  These four lectures sit outside Geron, so the notes ARE the primary
%  source and are examined on the same terms as the textbook parts.
%  Everything here is chosen to make that readable on paper: one column,
%  generous measure, and the same six-colour palette as the slides so a
%  student moving between the two is not re-learning what red means.
%
%  Build with pdflatex (twice, for the toc and the refs):
%      cd notes && latexmk -pdf lecture-19.tex
%  or  make -C notes
% =====================================================================
\usepackage[T1]{fontenc}
\usepackage[utf8]{inputenc}
\usepackage{newtxtext,newtxmath}      % Times-like: prints well, reads well
\usepackage[scaled=0.86]{inconsolata} % code, at a size that sits beside Times
\usepackage{microtype}
% newtx defines \Bbbk, \openbox, \proof and \endproof; amssymb and amsthm
% then redefine them and abort. Freeing the four names before those two
% packages load is the standard fix, and it costs nothing here: we use
% amsthm only for \newtheorem.
\let\Bbbk\relax
\usepackage{amsmath,amssymb}
\let\openbox\relax
\let\proof\relax
\let\endproof\relax
\usepackage{amsthm}
\usepackage{booktabs}
\usepackage{enumitem}
\usepackage{xcolor}
\usepackage{tcolorbox}
\usepackage{titlesec}
\usepackage{graphicx}
\usepackage[hidelinks]{hyperref}
\tcbuselibrary{skins,breakable}

% ---- the course palette, from assets/css/custom.css -------------------
\definecolor{aimlprimary}{HTML}{0B3D62}   % deep blue  - structure
\definecolor{aimlaccent} {HTML}{C0392B}   % brick red  - failures
\definecolor{aimlsuccess}{HTML}{14663A}   % green      - repairs
\definecolor{aimlmath}   {HTML}{6C3483}   % purple     - the derivation
\definecolor{aimlmuted}  {HTML}{4B5563}
\definecolor{aimlrule}   {HTML}{B0BCC7}
\definecolor{aimlsoft}   {HTML}{F4F7F9}

\usepackage[a4paper,margin=32mm,top=30mm,bottom=30mm]{geometry}
\setlength{\parskip}{0.45\baselineskip}
\setlength{\parindent}{0pt}
\linespread{1.06}

% ---- headings --------------------------------------------------------
\titleformat{\section}{\Large\bfseries\color{aimlprimary}}{\thesection}{0.7em}{}
\titleformat{\subsection}{\large\bfseries\color{aimlprimary}}{\thesubsection}{0.7em}{}
\titleformat{\subsubsection}{\bfseries\color{aimlmuted}}{}{0em}{}
\titlespacing*{\section}{0pt}{1.6\baselineskip}{0.5\baselineskip}
\titlespacing*{\subsection}{0pt}{1.2\baselineskip}{0.35\baselineskip}

% ---- boxes -----------------------------------------------------------
% One box per role, and the role is the colour. A student who has seen the
% slides already knows what each one means before reading a word of it.
\newtcolorbox{keybox}[1][]{
  breakable, enhanced, colback=aimlsoft, colframe=aimlprimary,
  boxrule=0.4pt, left=8pt, right=8pt, top=6pt, bottom=6pt,
  fonttitle=\bfseries\small, coltitle=white, colbacktitle=aimlprimary,
  attach boxed title to top left={yshift=-2mm, xshift=4mm},
  boxed title style={boxrule=0pt, arc=1pt}, #1}

\newtcolorbox{failbox}[1][]{
  breakable, enhanced, colback=white, colframe=aimlaccent,
  boxrule=0.9pt, left=8pt, right=8pt, top=6pt, bottom=6pt,
  fonttitle=\bfseries\small, coltitle=white, colbacktitle=aimlaccent,
  attach boxed title to top left={yshift=-2mm, xshift=4mm},
  boxed title style={boxrule=0pt, arc=1pt}, #1}

\newtcolorbox{derivbox}[1][]{
  breakable, enhanced, colback=white, colframe=aimlmath,
  boxrule=0.6pt, left=8pt, right=8pt, top=6pt, bottom=6pt,
  fonttitle=\bfseries\small, coltitle=white, colbacktitle=aimlmath,
  attach boxed title to top left={yshift=-2mm, xshift=4mm},
  boxed title style={boxrule=0pt, arc=1pt}, #1}

% An exercise. Deliberately the same shape as the notebook's `try` field:
% one change, and what should happen to the output.
\newtcolorbox{trybox}[1][]{
  breakable, enhanced, colback=white, colframe=aimlsuccess,
  boxrule=0.5pt, left=8pt, right=8pt, top=6pt, bottom=6pt,
  fonttitle=\bfseries\small, coltitle=white, colbacktitle=aimlsuccess,
  attach boxed title to top left={yshift=-2mm, xshift=4mm},
  boxed title style={boxrule=0pt, arc=1pt}, #1}

% ---- small semantics -------------------------------------------------
\newcommand{\scope}[1]{\textcolor{aimlmuted}{\small #1}}
\newcommand{\term}[1]{\textbf{#1}}
\newcommand{\code}[1]{\texttt{\small #1}}
\newcommand{\lecture}[1]{Lecture~#1}

\newtheorem{definition}{Definition}[section]
\newtheorem{proposition}[definition]{Proposition}

% ---- title block -----------------------------------------------------
% \notestitle{19}{Information retrieval: the lexical foundation}{SciFact (BEIR)}
\newcommand{\notestitle}[3]{%
  \thispagestyle{empty}
  \begin{flushleft}
    {\color{aimlmuted}\small\sffamily
     APPLICATIONS OF MACHINE LEARNING \quad$\cdot$\quad
     BSc Mathematics of Artificial Intelligence}\\[2mm]
    {\color{aimlrule}\rule{\textwidth}{0.8pt}}\\[4mm]
    {\color{aimlmuted}\large Lecture #1 \quad$\cdot$\quad Part V \quad$\cdot$\quad
     Extended notes}\\[3mm]
    {\Huge\bfseries\color{aimlprimary}\baselineskip=1.15\baselineskip #2\par}\vspace{4mm}
    {\color{aimlmuted}\large Dataset: #3}\\[3mm]
    {\color{aimlrule}\rule{\textwidth}{0.8pt}}
  \end{flushleft}
  \vspace{2mm}
  \begin{keybox}[title={Why these notes exist}]
    Lectures 19--22 sit outside G\'eron, the course's primary text. For those
    four lectures \emph{these notes are the primary source}, and they are
    examined on exactly the same terms as the textbook parts. Everything in
    the slide deck for this lecture is here, with the arguments written out at
    length rather than compressed onto a slide; the numbers are the ones the
    accompanying notebook prints, and every one of them is a measurement on
    the corpus named above rather than a figure quoted from a paper.
  \end{keybox}
  \vspace{1mm}
  {\small\tableofcontents}
  \vspace{2mm}
  \noindent{\color{aimlrule}\rule{\textwidth}{0.4pt}}
  \vspace{1mm}
}


\begin{document}
\notestitlefor{6}{I}{Decision trees}{Forest CoverType, seven species}{Chapter 5}

\section{A model family chosen by the shape of the answer}

Every model so far predicts with a weighted sum
$\boldsymbol\theta^{\intercal}\mathbf{x}$, fitted by solving a linear system or
descending a convex surface. Such models need their features scaled, because
the geometry of the loss depends on the units; they reach curves only by
manufacturing the curve yourself; and their explanation of a prediction is a
sum of $n$ terms --- a reason, but not a short one.

Today's model is not a weighted sum at all. It is a sequence of yes/no
questions, and \emph{the sequence is the explanation}.

\subsection{The brief}

An agency managing a national forest must publish a map of forest cover type
--- which species dominates each 30 by 30 metre patch --- over an area far too
large to survey on foot. From cartographic measurements alone (elevation,
slope, aspect, distances to water, road and fire ignition points, hillshade,
soil type) predict which of seven species covers the patch.

The map goes into the public record. It decides where logging is permitted,
which parcels qualify for habitat protection, and how fire risk is modelled for
the district. A landowner refused a permit is entitled to ask why. So is a
court.

\begin{keybox}[title={The requirement that decides the model}]
Every \emph{individual} prediction must be accompanied by a human-readable
justification: a statement, in terms of the measured quantities, of why this
patch was classified as that species.

Note the word \emph{individual}. Not ``which variables matter in general''. Not
``the model is 88\% accurate''. This patch, this answer, this reason.
\end{keybox}

``Human-readable'' is not a testable condition, so it has to be negotiated into
one before any code is written.

\begin{center}
\small
\begin{tabular}{p{0.30\textwidth}p{0.56\textwidth}}
\toprule
\textbf{Requirement} & \textbf{Testable form} \\
\midrule
Readable by a non-specialist & stated as conditions on the measured columns,
                               not on transformed ones \\
Short enough to check & at most 8 conditions per prediction \\
Auditable & the same patch gives the same justification, and the rule can be
            applied by hand \\
\bottomrule
\end{tabular}
\end{center}

Eight is the number the agency chose. It is arbitrary --- and it is now a
hyperparameter fixed by the brief rather than by cross-validation. That is
unusual and worth noticing.

\begin{center}
\small
\begin{tabular}{lp{0.52\textwidth}}
\toprule
\textbf{Model} & \textbf{Can it justify one prediction?} \\
\midrule
Logistic regression & partly --- a weighted sum of 54 terms is a reason, but
                      not a short one \\
$k$-nearest neighbours & ``because these 5 other patches looked similar'' \\
SVM with an RBF kernel & no --- the decision surface has no finite
                         description \\
Neural network & no \\
Decision tree & yes --- the path from root to leaf \emph{is} the
                justification \\
\bottomrule
\end{tabular}
\end{center}

This is the first problem in the course where the model family is chosen by the
shape of the answer rather than by the data.

\subsection{The data, and what we did to it}

581{,}012 patches, each with 54 measurements: 10 quantitative, 4 already
one-hot for wilderness area, 40 already one-hot for soil type. No missing
values, no text columns, nothing to impute. This one arrives clean, so the
difficulty is entirely elsewhere.

We use 60{,}000 patches, stratified so the class proportions are preserved
exactly, split 48{,}000 / 12{,}000.

\begin{failbox}[title={Failure condition --- numbers measured on a subsample}]
A 200-tree ensemble on all 581{,}012 rows takes minutes per fit, which does not
fit inside a ninety-minute session. Every number in this lecture and the next
is measured on the subsample, and every one of them would move on the full set.

Subsampling for speed is a legitimate choice and a \emph{reportable} one.
Subsampling and not mentioning it is not.
\end{failbox}

\begin{center}
\small
\begin{tabular}{lrl}
\toprule
\textbf{Class} & \textbf{Training patches} & \textbf{Species} \\
\midrule
2 & 23{,}405 & Lodgepole Pine \\
1 & 17{,}501 & Spruce/Fir \\
3 &  2{,}954 & Ponderosa Pine \\
7 &  1{,}694 & Krummholz \\
6 &  1{,}435 & Douglas-fir \\
5 &    784 & Aspen \\
4 &    227 & Cottonwood/Willow \\
\bottomrule
\end{tabular}
\end{center}

The commonest class is a hundred times the rarest, and two species are 85\% of
the ground. A classifier that always says ``Lodgepole Pine'' scores
\textbf{48.8\%}, so anything below that is worse than a constant and the
distance above it is the only part you earned. The agency's surveyors are
stated to agree with one another on roughly nine patches in ten --- an
assumption we were handed, not a measurement we made --- so the useful range
lies between those two numbers.

\section{Derivation: impurity, and the objective CART minimises}

You called \code{DecisionTreeClassifier().fit()}. It chose one feature and one
threshold at every node. Chose them to minimise \emph{what}?

A node holds $m$ instances over $K$ classes, with $p_k$ the fraction in class
$k$. The class distribution $\mathbf{p}$ is the \emph{only} thing an impurity
measure sees. A node is \term{pure} when some $p_k = 1$, and \term{uniform}
when every $p_k = 1/K$.

\subsection{What an impurity measure has to do}

Before choosing a formula, say what any acceptable one must satisfy:
\[
\begin{aligned}
\text{(i)}\;\; & I(\mathbf{p}) = 0 \iff \mathbf{p}\ \text{is pure} \\
\text{(ii)}\;\; & I\ \text{is maximal at}\ \mathbf{p} = (1/K,\dots,1/K).
\end{aligned}
\]
A node we have finished with scores zero; the node we know least about scores
worst. Everything else is a choice \emph{within} this specification, not a
disagreement about it.

Both of Scikit-Learn's criteria have one shape:
\[ I(\mathbf{p}) = \sum_{k=1}^{K}\varphi(p_k), \]
with $\varphi:[0,1]\to\mathbb{R}$ strictly concave and
$\varphi(0)=\varphi(1)=0$. Concavity is not decoration --- it is the property
that delivers both conditions above and the guarantee below, and it is the only
property either proof uses.

\term{Gini impurity} takes $\varphi(p) = p(1-p)$:
\[ G(\mathbf{p}) = \sum_k p_k(1-p_k) = 1 - \sum_k p_k^{\,2}. \]
Draw two of the node's instances independently with replacement;
$\sum_k p_k^2$ is the probability they are the same class, so $G$ is the
probability they \emph{differ}. A pure node never disagrees with itself.

\term{Entropy} takes $\varphi(p) = -p\log_2 p$:
\[ H(\mathbf{p}) = -\sum_k p_k\log_2 p_k, \]
the average number of bits needed to transmit the class of an instance drawn
from the node, under the best possible code. A pure node needs none.

\begin{derivbox}[title={Both satisfy the specification}]
\emph{(i)} $\varphi \ge 0$ on $[0,1]$ and vanishes only at $p \in \{0,1\}$, for
both choices. So $I(\mathbf{p}) = 0$ forces every $p_k \in \{0,1\}$, and since
they sum to 1 exactly one is 1: the node is pure.

\emph{(ii)} By Jensen's inequality on the concave $\varphi$,
\[ \frac{1}{K}\sum_k\varphi(p_k)
   \le \varphi\Bigl(\frac{1}{K}\sum_k p_k\Bigr) = \varphi(1/K), \]
so $I(\mathbf{p}) \le K\varphi(1/K)$. Strict concavity makes this an equality
only when every $p_k$ is equal, so the maximum is attained at the uniform
distribution and nowhere else.
\end{derivbox}

\begin{center}
\small
\begin{tabular}{rll}
\toprule
$K$ & $G_{\max} = 1 - 1/K$ & $H_{\max} = \log_2 K$ (bits) \\
\midrule
2 & 0.5 & 1 \\
7 & $6/7 \approx 0.857$ & $\log_2 7 \approx 2.807$ \\
\bottomrule
\end{tabular}
\end{center}

The two are on \emph{different scales}. Any comparison of their shapes has to
divide one of them first, and saying which is part of the comparison.

\subsection{The objective, and why a split never hurts}

A split sends $m_{\text{L}}$ instances left and $m_{\text{R}}$ right. Its cost
is the impurity of the two children, weighted by their size:
\[ J(k, t_k) = \frac{m_{\text{L}}}{m}I_{\text{L}}
             + \frac{m_{\text{R}}}{m}I_{\text{R}}, \]
searched over every feature $k$ and every threshold $t_k$. Equivalently,
maximise $\Delta I = I_{\text{parent}} - J$; the parent term is the same for
every candidate at this node, so the two are the same search.

\begin{derivbox}[title={$\Delta I \ge 0$, always}]
Every instance goes left or right, so for each class the parent share is the
size-weighted mixture of the children's shares:
\[ p_k = \frac{m_{\text{L}}}{m}p_{\text{L},k}
       + \frac{m_{\text{R}}}{m}p_{\text{R},k}. \]
Apply concavity of $\varphi$ to that convex combination and sum over $k$:
\[ I(\mathbf{p}) \ge \frac{m_{\text{L}}}{m}I_{\text{L}}
                   + \frac{m_{\text{R}}}{m}I_{\text{R}} = J. \]
\end{derivbox}

So the greedy search can always stop safely: no split ever makes things worse,
and $\Delta I = 0$ signals that this node has nothing left to gain.

\subsection{Why the two criteria usually choose the same split}

$I(\mathbf{p}) = \sum_k\varphi(p_k)$ with $\varphi$ concave is
\term{Schur-concave}: if one distribution majorises another --- is more
concentrated, in the sense that its largest share is larger, its two largest
shares sum to more, and so on --- then it has the lower impurity, \emph{for
every such $\varphi$}.

So take two candidate splits sending the same numbers left and right. If A's
two children each majorise B's, then $J_{\text{A}} \le J_{\text{B}}$ for every
concave $\varphi$, and Gini and entropy must agree that A is better. Neither
the choice of $\varphi$ nor its scale can change it.

They are free to disagree only on pairs that are \emph{not} comparable in that
ordering. How far can they part? On a two-class node, $G$ has maximum $1/2$ and
$H$ has maximum 1, so compare $H/2$ with $G$: both are 0 at $p \in \{0,1\}$ and
$1/2$ at $p = 1/2$, and the largest gap anywhere else is \textbf{0.055}, at
$p \approx 0.10$ and again at $p \approx 0.90$. Two candidates whose weighted
impurities differ by more than 0.055 are ranked identically by both criteria.
Disagreement needs a close call \emph{and} an incomparable pair.

\subsection{Measure it rather than accepting it}

Géron says entropy tends to produce slightly more balanced trees, that Gini is
marginally faster, and that ``most of the time it does not make a big
difference''. That last is a claim about behaviour, and we have a dataset.
Twenty resamples, both criteria on the same resample each time, so the
comparison is paired:

\begin{center}
\small
\begin{tabular}{lr}
\toprule
\textbf{Question} & \textbf{Measured} \\
\midrule
Resamples where the two chose the same root feature & 20 of 20 \\
Test predictions on which the two trees agree & 85.3\% \\
Mean accuracy, gini & 73.4\% \\
Mean accuracy, entropy & 73.0\% \\
Paired difference, entropy $-$ gini & $-0.44 \pm 0.39$ points \\
\bottomrule
\end{tabular}
\end{center}

Exactly the shape the Schur-concavity argument predicts: the two pick the same
first question every time, because the root split is not a close call, and then
disagree about one prediction in seven, deep in the tree where candidates are
separated by a hair.

\begin{keybox}[title={Three true statements, and only the third is a recommendation}]
The effect is \emph{real}: Gini wins on 17 of the 20 resamples and the mean
sits about five standard errors from zero. The effect is \emph{tiny}: 0.44
points, against the 8.0 points that \code{max\_depth} is worth later in this
lecture. So do not spend your tuning budget here.

``Statistically detectable'' and ``worth acting on'' are different claims and
need different evidence. Twenty paired resamples give you both; one run gives
you neither.
\end{keybox}

\subsection{What the derivation does not give you}

$\Delta I \ge 0$ guarantees each split is locally harmless. It says nothing
about the \emph{sequence} of splits, and CART chooses that greedily: the best
split at this node, never reconsidered. A split that looks mediocre now but
would enable two excellent splits below it is never found.

Finding the optimal tree is NP-complete; the greedy algorithm returns a good
one in $O(n\,m\log m)$ and no known algorithm returns the best one in
reasonable time. So the tree you get is a function of the algorithm as much as
of the data --- the seed of this lecture's second failure condition.

\section{Growing a tree, and reading one}

Each internal node asks $x_k \le t_k$ about one feature; each leaf predicts the
majority class of the training instances that reached it. Predicting one
instance costs $O(\log_2 m)$ --- one root-to-leaf path --- and does not depend
on $n$, which is why reading a prediction costs the same as making it. That is
what makes the justification free.

Trees need no feature scaling at all: $x_k \le t_k$ is invariant under any
strictly increasing transformation of $x_k$, so standardising a column changes
its thresholds and nothing else.

\subsection{Fit one with no constraints}

Depth 33, 5{,}699 leaves, 100.0\% training accuracy. With 48{,}000 training
patches that is eight patches per leaf on average, so of course it is perfect
on its own data.

\begin{center}
\small
\begin{tabular}{lr}
\toprule
\textbf{Unconstrained tree} & \textbf{Accuracy} \\
\midrule
On the training set & 100.0\% \\
5-fold cross-validated, on the training set & 82.1\% \\
On the test set (once, at the end) & 82.6\% \\
\bottomrule
\end{tabular}
\end{center}

82\% is real and far above the 48.8\% baseline, so the tree family works here.
But count what a prediction rests on. Using \code{decision\_path}, the mean
justification is \textbf{17.88} conditions and the worst is \textbf{33}. The
brief allows eight. (The maximum equalling \code{get\_depth()} is the cheap
check that the off-by-one went the right way.)

It grew to depth 33 because nothing stopped it. A tree is a \term{nonparametric}
model: the parameter count is not fixed before training, so the structure
follows the data as far as the data goes. A linear model cannot do this --- its
parameter count is decided by the number of columns before it sees a row.

\subsection{What depth actually buys}

\begin{center}
\small
\begin{tabular}{rrrl}
\toprule
\code{max\_depth} & \textbf{CV accuracy} & \textbf{Leaves} &
\textbf{What a justification looks like} \\
\midrule
1  & 63.3\% & 2 & one condition \\
3  & 67.5\% & 8 & a sentence \\
8  & 73.8\% & 206 & a short paragraph --- the brief \\
18 & 81.7\% & 3{,}377 & nobody reads this \\
\bottomrule
\end{tabular}
\end{center}

Between depth 8 and depth 18 there are \textbf{8.0 points} of accuracy. That is
the measured price of the eight-condition requirement, and it is a number worth
putting in front of the agency.

\begin{failbox}[title={Failure condition --- the one hyperparameter we may not tune}]
Cross-validation says \code{max\_depth=None}. The brief says
\code{max\_depth=8}. Cross-validation does not get a vote, because it is not
optimising the thing the agency is buying.

This is the first time in the course that the honest answer to ``what does the
search say?'' is ``the search is answering the wrong question''. It will not be
the last. Bring the 8.0-point figure to the agency; perhaps they will trade a
condition or two for it. That is a conversation, not a \code{GridSearchCV}.
\end{failbox}

We search the whole grid anyway --- including the depths we may not use ---
because the rows we cannot pick are exactly what tells the agency what its
requirement costs. Two things fall out. Within the permitted depth-8 row
\code{min\_samples\_leaf} barely matters, 73.8\% at 1 against 71.5\% at 200,
because the depth limit binds first. Down the unlimited-depth row it matters
enormously, 82.1\% at 1 falling to 71.9\% at 200, because there the leaf size
\emph{is} the regularisation. Two hyperparameters that both restrict the tree
do not act independently, and a 2-D grid shows that where two separate 1-D
sweeps would not.

\subsection{Overruling the grid a second time}

The grid's answer under the cap is \code{min\_samples\_leaf=1}. We ship 20
anyway.

\begin{keybox}[title={Why}]
The model states its justification as ``91\% of the 463 training patches in
this leaf''. At \code{min\_samples\_leaf=1} that becomes ``100\% of the 1'' ---
a single surveyed patch wearing the grammar of evidence. The brief requires a
justification that can be audited; that is not one.

Cross-validation prefers 1 at 73.75\%; the brief requires 20 at 73.35\%. The
price of auditability is 0.40 points, and it goes to the agency. A constraint
that overrules a measurement has to show the measurement it overruled.
\end{keybox}

\begin{center}
\small
\begin{tabular}{lrr}
\toprule
& \textbf{Unconstrained} & \textbf{Depth 8, leaf 20} \\
\midrule
Leaves & 5{,}699 & 163 \\
Columns consulted, of 54 & 47 & 24 \\
Mean conditions per justification & 17.88 & 7.97 \\
Smallest leaf & 1 & 20 \\
Training accuracy & 100.0\% & 74.4\% \\
Cross-validated accuracy & 82.1\% & 73.4\% \\
\bottomrule
\end{tabular}
\end{center}

Notice what the last two rows now do: 74.4\% against 73.4\%, a gap of \emph{one
point}, where the unconstrained tree's gap was 17.9. The depth limit did not
only shorten the justification --- it removed almost all of the overfitting as
a side effect.

Which leaves a question worth holding on to. If the constrained tree barely
overfits, why is it still nearly nine points worse than the unconstrained one?

\subsection{Reading one prediction}

The entire justification mechanism is \code{tree\_.children\_left},
\code{tree\_.feature} and \code{tree\_.threshold} --- nine lines, no library
beyond what is already imported. For one test patch:

\begin{center}
\small
\begin{verbatim}
1. Elevation   = 3,763 >  3,046
2. Elevation   = 3,763 >  3,306
3. Wild_2      =     1 >      0
4. Soil_31     =     0 <=     0
5. Elevation   = 3,763 >  3,360
6. HDist_Fire  = 3,020 >    364
7. Shade_Noon  =   200 <=   244
8. Shade_3pm   =   133 <=   196
-> Krummholz  (91% of the 463 training patches in this leaf)
\end{verbatim}
\end{center}

Eight conditions, every one a measurement a surveyor took. The path \emph{is}
the justification: no extra machinery, no approximation, no second model
explaining the first.

A note on thresholds: CART puts the split midway between two adjacent observed
values, which is why the root reads $3046.5$ rather than $3046$. Nothing in the
data sits at 3{,}046.5 metres.

\begin{failbox}[title={Failure condition --- read the leaf carefully}]
91\% of the 463 training patches in that leaf were Krummholz. That is a
statement about 463 \emph{training patches}, not a probability that this patch
is Krummholz.

\code{predict\_proba} returns exactly those leaf proportions, so a tree's
``probabilities'' are piecewise constant and identical for every patch reaching
the same leaf. Across the 163 leaves the smallest holds 20 patches and the
largest several thousand; those deserve very different amounts of trust, and
the count belongs in the justification.

The honest sentence is ``of the training patches that satisfied these eight
conditions, 91\% were Krummholz'' --- not ``this patch is 91\% likely to be
Krummholz''.
\end{failbox}

\section{One number for seven species is a summary}

The test set, once: \textbf{73.0\%} on 12{,}000 patches, against a constant
baseline of 48.8\%, with every choice made on cross-validated folds of the
training set.

\begin{center}
\small
\begin{tabular}{lrr}
\toprule
\textbf{Cover type} & \textbf{Share of the data} & \textbf{Recall} \\
\midrule
Lodgepole Pine     & 48.8\% & 82.0\% \\
Ponderosa Pine     &  6.2\% & 80.8\% \\
Spruce/Fir         & 36.5\% & 69.6\% \\
Krummholz          &  3.5\% & 56.4\% \\
Cottonwood/Willow  &  0.5\% & 56.1\% \\
Douglas-fir        &  3.0\% & 12.8\% \\
Aspen              &  1.6\% &  2.6\% \\
\bottomrule
\end{tabular}
\end{center}

The model finds fewer than 3 Aspen patches in every hundred. On the headline
number that costs 1.6 points and is invisible.

\begin{failbox}[title={Failure condition --- why the tree cannot see Aspen}]
$J$ weights each child's impurity by \emph{how many instances it holds}. A
split that isolates a class at 1.6\% of the node moves the weighted average by
almost nothing, so the greedy search never chooses one. With
\code{max\_depth=8} there are at most 256 leaves to allocate across seven
species, and the algorithm spends them where the instances are.

This is Lecture 3's lesson in a new costume: an objective that averages over
instances rewards ignoring the rare class. It is a property of the objective,
not a bug in the implementation --- read the definition of $J$ again and you
could have predicted it.
\end{failbox}

\code{class\_weight="balanced"} multiplies each instance's contribution by the
inverse of its class frequency. Aspen recall rises to most of the class,
because the split that isolates it is now worth choosing; overall accuracy
falls sharply, because the two species that are 85\% of the ground now carry
$2/7$ of the weight instead of 85\% of it. A different model answering a
different question --- and quoting either number without the other is an
argument rather than a measurement.

\begin{keybox}[title={The sentence the agency needs}]
``This map is 73.0\% accurate overall, but it identifies fewer than three Aspen
stands in a hundred. Do not use it to decide anything about Aspen or
Douglas-fir.''

A model can be fit for one purpose and unfit for another in the same
deployment. Saying which is which is the job.
\end{keybox}

\section{Regression trees, and two ways a tree breaks}

\code{DecisionTreeRegressor} is the same algorithm with the impurity replaced
by the mean squared error of the node about its own mean. A leaf predicts the
mean of the training targets that reached it, so the output is a step function,
one step per leaf. An unconstrained regression tree drives its training error
to exactly zero. And a tree \emph{cannot extrapolate}: beyond the range of the
training targets it returns the nearest leaf's mean forever, which is the thing
to know before using one on a trending series.

\subsection{Failure condition 1 --- every boundary is axis-aligned}

Every split is $x_k \le t_k$, so every decision boundary is perpendicular to an
axis. A boundary that genuinely runs at $45^\circ$ has to be approximated by a
staircase, and each step costs a level of depth. Rotate the data by
$45^\circ$ and the same problem needs a much deeper tree, though nothing about
the problem changed.

Scaling a column is harmless --- it is monotone in that column alone. Rotating
the space is not: it mixes columns, and the family of candidate splits is not
rotation-invariant.

\subsection{Failure condition 2 --- ask a question accuracy cannot answer}

You have a model whose every prediction comes with a rule a person can read.
Fit it again --- same hyperparameters, same seed --- on 90\% of the same
training set. How similar are the two sets of rules? Almost everybody expects
``very''. The answer has two halves and they point in opposite directions.

\begin{center}
\small
\begin{tabular}{lr}
\toprule
\textbf{Over 20 refits} & \textbf{Measured} \\
\midrule
Accuracy, mean & 72.99\% \\
Accuracy, standard deviation & 0.25\% \\
Accuracy, range & 72.39\% -- 73.38\% \\
Pairwise prediction disagreement, over all 190 pairs & 9.1\% \\
\quad worst pair & 15.4\% \\
Test patches all twenty trees agree on & 71.4\% \\
\bottomrule
\end{tabular}
\end{center}

Two trees with the same accuracy disagree about the species of one patch in
eleven. And where in the tree does it change?

\begin{center}
\small
\begin{tabular}{lr}
\toprule
\textbf{Top of the tree} & \textbf{Across the 20 refits} \\
\midrule
Feature at the root & Elevation, 20 times out of 20 \\
Distinct thresholds at the root & 4, spanning 15 metres of elevation \\
Distinct structures in the top three levels & 11 of 20 \\
Leaves in the whole tree & 151 to 170 \\
Columns the tree consults & 22 to 27 of 54 \\
\bottomrule
\end{tabular}
\end{center}

\begin{keybox}[title={The part you put on a slide is stable. The part that decides is not.}]
A tree is a hierarchy. The root split is chosen from 48{,}000 patches and wins
by a wide margin, so it is stable --- which is also why Gini and entropy agreed
about it 20 times out of 20. A node eight levels down was chosen from a few
hundred patches, where two candidates are often separated by a hair; change one
row, the loser becomes the winner, and everything below is a different tree.

Accuracy is an average over 12{,}000 patches, and averages hide substitutions:
a tree that loses a patch here and gains one there scores the same. \emph{A
stable metric is not evidence of a stable model. It is evidence that you
measured the wrong thing.}
\end{keybox}

This is \emph{variance} in the sense of Lecture 5's decomposition --- the
sensitivity of the fitted model to the particular training sample. Trees have
it badly for a structural reason: the greedy, hierarchical fit turns an
arbitrarily small change in one split into an arbitrarily large change in the
subtree below it, and nothing smooths it out.

Failure condition 1 was the same fact in a different hat. Rotating the data
changes which axis-aligned split wins by a hair, and the tree below it changes
completely.

The next lecture derives what averaging does to variance and builds the models
that exploit it. It also shows what the repair costs: the justification we
spent this lecture earning.

\section{The notebook}

\begin{trybox}[title={What to change}]
Refit the shipped tree on 90\% subsamples with a different set of seeds and
recompute the pairwise disagreement. If it lands near 9\% again, you have
reproduced the instability rather than read about it --- and you can then ask
the harder question: which of the 12{,}000 patches are the unstable ones, and
do they have anything in common?
\end{trybox}

\section{Exercises}

\begin{enumerate}[leftmargin=1.6em,itemsep=4pt]
  \item State the two conditions any impurity measure must satisfy, and prove
        that entropy attains its maximum only at the uniform distribution.
        \hfill \scope{[5]}
  \item Show that a split can never increase the size-weighted impurity, naming
        the property of $\varphi$ the argument uses. \hfill \scope{[4]}
  \item Two candidate splits send the same numbers left and right, and each of
        A's children majorises the corresponding child of B. State which
        criterion prefers A, and why the answer does not depend on the choice
        between Gini and entropy. \hfill \scope{[4]}
  \item A tree achieves 73.0\% accuracy but 2.6\% recall on a class holding
        1.6\% of the data. Explain, from the definition of the split
        objective, why this is expected rather than a defect.
        \hfill \scope{[5]}
  \item Twenty trees fitted on 90\% subsamples of one training set have
        accuracies within 0.5 points of each other but disagree on 9\% of test
        predictions. Explain how both can be true, and say what it tells you
        about reporting accuracy alone. \hfill \scope{[4]}
\end{enumerate}

\section*{Summary}
\addcontentsline{toc}{section}{Summary}

The model family was chosen by the shape of the required answer rather than by
the data: only a tree can hand a court an individual justification in the
measured columns. Impurity is $\sum_k\varphi(p_k)$ for strictly concave
$\varphi$, which is the only property the theory uses --- it gives zero exactly
at purity, a maximum exactly at uniformity by Jensen, and $\Delta I \ge 0$ for
every split by the same concavity. Schur-concavity explains why Gini and
entropy agree wherever the comparison is not a close call, and the measurement
agreed: the same root feature 20 times out of 20, and 0.44 points between them
against the 8.0 that depth is worth.

The brief fixed \code{max\_depth=8}, which cost 8.0 points, and auditability
fixed \code{min\_samples\_leaf=20}, which cost 0.40 more. Both figures go to
the agency, because a constraint that overrules a measurement has to show the
measurement it overruled. The shipped tree scores 73.0\% against a 48.8\%
constant, with a mean justification of 7.97 conditions.

Then two failure conditions, both properties of the algorithm. Every boundary
is axis-aligned, so a $45^\circ$ boundary costs depth. And twenty refits on
90\% subsamples held accuracy to within a quarter of a point while disagreeing
on 9.1\% of predictions --- variance, in Lecture 5's sense, arriving as a
stable metric hiding an unstable model.

\end{document}
